\chapter{Проверка соответствия требованиям} \label{ch5}

\section{Требования к раздаточному материалу} \label{ch5:report}

\noindent
\begingroup
\centering
\small
\begin{longtable}[c]{|p{6.5cm}|p{8.5cm}|}
    \caption{Соответствие отчета обязательным элементам}
    \label{tab:report-requirements}
    \\
    \hline
    Требование & Реализация в работе \\ \hline
    \endfirsthead
    \hline
    Требование & Реализация в работе \\ \hline
    \endhead
    \hline
    Тема относится к автоматизации или теории тестирования & Исследуется автоматизированный запуск модульных тестов в Jest, Vitest и Mocha. \\ \hline
    Использование принципа IMRAD & Введение, методы, результаты, обсуждение и заключение выделены структурно. \\ \hline
    Математическая постановка проблемы & Раздел~\ref{ch1:math}. \\ \hline
    Иллюстративный пример & Раздел~\ref{ch1:example}. \\ \hline
    Псевдокод алгоритма в окружении \texttt{algorithm} & Раздел~\ref{ch1:algorithm}, алгоритм~\ref{alg:benchmark}. \\ \hline
    Пошаговое выполнение алгоритма на данных & Раздел~\ref{ch1:steps}. \\ \hline
    Описание разработанного программного обеспечения & Глава~\ref{ch2}. \\ \hline
    Подробное описание средств и подходов тестирования & Разделы~\ref{ch1:frameworks}, \ref{ch2:testing}. \\ \hline
    Инструкция по установке и запуску & Раздел~\ref{ch2:run}. \\ \hline
    Автоматически сгенерированный список источников из bib-файла & Используется \texttt{my\_folder/my\_biblio.bib}; список выводится командой \texttt{\textbackslash printbibliography}. \\ \hline
\end{longtable}
\normalsize
\endgroup

\section{Требования к программному обеспечению} \label{ch5:software}

\noindent
\begingroup
\centering
\small
\begin{longtable}[c]{|p{6.5cm}|p{8.5cm}|}
    \caption{Соответствие программы требованиям задания}
    \label{tab:software-requirements}
    \\
    \hline
    Требование & Реализация \\ \hline
    \endfirsthead
    \hline
    Требование & Реализация \\ \hline
    \endhead
    \hline
    Комментирование кода & В коде добавлены комментарии к детерминированному расширению данных, сэмплированию процесса и очистке кеша pidusage. \\ \hline
    Компактность кода и проекта & Основная логика находится в трех файлах каталога \texttt{src}. \\ \hline
    Быстрое разворачивание и компиляция при code review & Достаточно выполнить \texttt{npm install}, затем \texttt{npm run verify}. \\ \hline
    Наличие \texttt{requirements.txt} & Файл находится в корне \texttt{benchmark\_app}. \\ \hline
    Входные данные переменной величины & JSON-поле \texttt{size} задает итоговую длину выборки. \\ \hline
    Импорт CSV или JSON & Реализована загрузка обоих форматов. \\ \hline
    Эталонный входной файл и эталонный выход & \texttt{reference-input.json} и \texttt{reference-output.json}. \\ \hline
    Авторский иллюстративный пример & \texttt{author-example.json} и \texttt{author-example.csv}. \\ \hline
\end{longtable}
\normalsize
\endgroup

\section{Файлы результатов} \label{ch5:reports}

После запуска эксперимента формируются четыре файла:
\begin{enumerate}
    \item \texttt{reports/benchmark-results-json.json};
    \item \texttt{reports/benchmark-results-json.csv};
    \item \texttt{reports/benchmark-results-csv.json};
    \item \texttt{reports/benchmark-results-csv.csv}.
\end{enumerate}
JSON-файлы подходят для машинной обработки, CSV-файлы можно открыть в электронных таблицах.
